\pdfinfoomitdate=1
\pdftrailerid{}
\pdfsuppressptexinfo=15
\newcommand{\DoNotLoadEpstopdf}{}
\documentclass[11pt]{article}
\usepackage[T1]{fontenc}
\usepackage[utf8]{inputenc}
\usepackage{lmodern,amsmath,amssymb,amsthm,mathtools,booktabs,array}
\usepackage[margin=1in]{geometry}
\usepackage{microtype}
\usepackage[hidelinks]{hyperref}
\usepackage{enumitem}
\setlist{itemsep=3pt,topsep=5pt}
\allowdisplaybreaks[2]
\emergencystretch=2em
\newtheorem{theorem}{Theorem}[section]
\newtheorem{lemma}[theorem]{Lemma}
\newtheorem{proposition}[theorem]{Proposition}
\newtheorem{corollary}[theorem]{Corollary}
\theoremstyle{definition}\newtheorem{definition}[theorem]{Definition}
\theoremstyle{remark}\newtheorem{remark}[theorem]{Remark}
\newcommand{\NN}{\mathbb N_0}
\newcommand{\Id}{\mathrm{Id}}
\newcommand{\OO}{\mathcal O}
\newcommand{\PP}{\mathcal P}
\newcommand{\norm}[1]{\left\lVert#1\right\rVert}
\newcommand{\rhol}{\rho_{\mathrm{lo}}}
\newcommand{\rhoh}{\rho_{\mathrm{hi}}}
\newcommand{\cl}{c_{\mathrm{lo}}}
\newcommand{\ch}{c_{\mathrm{hi}}}
\DeclareMathOperator{\Res}{Res}
\DeclareMathOperator{\ord}{ord}
% Added by the ProveIt write phase (batch 110, 7 October 2026): dated notes
% marked [write], a macro for file, label and declaration names, longtable.
\usepackage{longtable}
\newcommand{\file}[1]{\nolinkurl{#1}}
\expandafter\def\expandafter\UrlBreaks\expandafter{\UrlBreaks\do\-\do\_\do\.}
\newenvironment{writenote}[1][]{\par\smallskip\begingroup\small
 \noindent\textbf{[write]}\if\relax\detokenize{#1}\relax\else~\textbf{#1.}\fi\ }%
 {\par\endgroup\smallskip}
\hypersetup{pdftitle={Report183 Certified asymptotics for tableau defect clusters A394326},pdfauthor={Research report},pdfsubject={Signed Fredholm determinants, certified dominant pole, spectral expansions and threshold inverses}}
\title{Report183\\[5pt]Certified asymptotics for tableau defect clusters\\[3pt]OEIS A394326}
\author{A mathematical and computational reproducibility report}
\date{3 October 2026}
\begin{document}
\maketitle
\begin{abstract}
The connected inversion-deficit cluster series associated with rectangular three-row standard Young tableaux has an exponential growth constant close to, but strictly smaller than, the golden ratio. We identify the source's hard-rod extraction with an exact first-return quotient and construct a signed trace-class operator whose Fredholm determinant controls its poles. An integer-only rational-circle certificate, outward dyadic local calculations, and explicit infinite-tail bounds prove a unique noncancelled simple dominant pole. With $a_0=0$, the resulting uniform estimate is
\[
 \left|a_n-c\rho^{-n}\right|\le 2500(100/63)^n\qquad(n\ge0),
\]
where $0.6180397469174016<\rho<0.6180397783636036$ and
$0.1838102868<c<0.1838120734$. Meromorphic continuation throughout the open unit disk gives every fixed finite spectral order, including multiple and complex poles. We prove eventual strict growth, explicit two-ceiling threshold-inverse bounds, all-order Lagrange expansions of their bounding roots, and a local signed span-marker theorem. The source's strict golden-ratio and Fibonacci equivalents are false for this series; its weaker small-residual observation remains compatible with the result.
\end{abstract}

\paragraph{Provenance and scope.}
The sequence and its extraction are due to Morten Brydensholt in the OEIS record; Sean A. Irvine's jOEIS translation supplies the inspected executable specification~\cite{oeis,joeis}. The report gives target-specific proofs and reproducible certificates. The methods of ballot paths, renewal, Fredholm determinants, complex coefficient extraction, and Lagrange inversion are classical. The bounded source screen described in Section~\ref{tdc:sec:sources} is not a worldwide priority determination. No publication or external record edit is part of this work.

\begin{writenote}[Status and trust boundary, added 7 October 2026, batch~110]
This is bundle Report~183 of an external AI-assisted research session, filed
in ProveIt as the single-source research report
\file{a394326-tableau-defect-clusters} (provenance, the sources as the write
read them, what was checked, relation to the repository, the collected
non-claims and the reading conventions are in
Section~\ref{tdc:sec:provenance}). It is unrefereed and not formalized. Its
author line reads ``A mathematical and computational reproducibility
report'' and its PDF author field ``Research report'': it names no person or
tool. \emph{What rests on what.} The identification of the source extraction
(Section~\ref{tdc:sec:model}), the operator construction and the analytic
arguments are conventional proofs. Theorem~\ref{tdc:thm:main} also rests on
a computer certificate: an exact integer rational-circle computation
(Section~\ref{tdc:sec:global}) and outward dyadic interval linear algebra
with frozen proposals (Section~\ref{tdc:sec:local}), replayed by the shipped
standard-library programs; these are not proof-assistant checked.
Corollary~\ref{tdc:cor:empirical}, Theorem~\ref{tdc:thm:inverse} and the
monotone threshold $Y_0$ inherit that trust boundary.
Theorems~\ref{tdc:thm:spectral} and~\ref{tdc:thm:marker} are existential
beyond the dominant pole. \emph{A correction of the OEIS record.} The two
asymptotic formulas of A394326, $a(n)\sim A\phi^n$ and
$a(n)\sim a(n-1)+a(n-2)$, are false; the source proves this
(Corollary~\ref{tdc:cor:empirical}), and the write records it with the
quotations and an explicit proof in Remark~\ref{tdc:rem:refuted}. The write
also adds Remark~\ref{tdc:rem:oeis} (the OEIS entries quoted),
Remark~\ref{tdc:rem:transseries} (the inverses against the transseries
volume) and dated notes. No statement, proof or number of the source was
changed.
\end{writenote}

\tableofcontents
\newpage

\section{Results and how to use them}
\label{tdc:sec:main}
Let $a_n$ be A394326 for $n\ge1$, extend it by $a_0=0$, and put
\[
 C(q)=\sum_{n\ge0}a_nq^n,\qquad R=\frac{63}{100},\qquad M=2500.
\]
The first terms, starting at $n=1$, are
\[
 1,2,1,1,3,5,6,7,15,27,39,57,97,161,257,405,659,1076.
\]
Section~\ref{tdc:sec:model} gives a complete definition independent of this prefix. In particular, the later asymptotic statements apply to the exact source extraction, not to a fitted recurrence.

\begin{theorem}[Certified dominant pole]\label{tdc:thm:main}
The series $C$ has a meromorphic continuation to $|q|<1$. There exist real constants $\rho,c$ with
\begin{align}
 \rhol:=0.6180397469174016&<\rho<0.6180397783636036=:\rhoh,
 \label{tdc:eq:rho}\\
 \cl:=0.1838102868&<c<0.1838120734=:\ch
 \label{tdc:eq:c}
\end{align}
such that $q=\rho$ is the only singularity of $C$ in $|q|\le R$. It is a simple pole, with principal term $c/(1-q/\rho)$. For every integer $n\ge0$,
\begin{equation}\label{tdc:eq:error}
 \left|a_n-c\rho^{-n}\right|\le MR^{-n}.
\end{equation}
Equivalently, $a_n=c\rho^{-n}\{1+O((\rho/R)^n)\}$. The growth constant $\lambda=\rho^{-1}$ satisfies
\begin{equation}\label{tdc:eq:lambda}
 1.6180188314<\lambda<1.6180189139<\frac{1+\sqrt5}{2}.
\end{equation}
The constants are defined exactly by the operators in Section~\ref{tdc:sec:operator}:
\begin{equation}\label{tdc:eq:exactconstants}
 D(q)=\det(\Id-K(\sqrt q)),\quad N(q)=\det(\Id-L(\sqrt q)),
 \quad N(\rho)=0,\quad c=\frac{D(\rho)}{\rho N'(\rho)}.
\end{equation}
There is exactly one zero of $N$ in $|q|<R$, counted with multiplicity, and $N$ has no boundary zero. This zero is $\rho$ and $D(\rho)\ne0$.
\end{theorem}

All decimal endpoints in \eqref{tdc:eq:rho}--\eqref{tdc:eq:c} denote exact terminating decimals. They are deliberately wider than the raw dyadic enclosures. Numerical refinements
\[
 \rho\approx0.6180397634035489,\qquad c\approx0.183811179648135
\]
are useful diagnostics only; the extra displayed digits are not certified here.

\begin{corollary}[Growth and the empirical equivalents]\label{tdc:cor:empirical}
For every $n\ge600$, $a_{n+1}>a_n>0$. If $\phi=(1+\sqrt5)/2$, then
\begin{equation}\label{tdc:eq:empirical}
 \frac{a_n}{\phi^n}\longrightarrow0,\qquad
 \frac{a_n-a_{n-1}-a_{n-2}}{a_n}\longrightarrow1-\rho-\rho^2
 \in(-0.00001295,-0.00001287).
\end{equation}
Thus neither a nonzero golden-ratio equivalent nor the strict equivalent
$a_n\sim a_{n-1}+a_{n-2}$ holds. The source's looser observed absolute relative residual below $0.0002$ is not contradicted.
\end{corollary}

The stronger structural conclusions have different scopes. Theorem~\ref{tdc:thm:spectral} supplies a finite-pole expansion at every fixed zero-free radius below one; it does not numerically locate later poles. Theorem~\ref{tdc:thm:inverse} gives an explicit large-threshold inverse enclosure. Its two real bounding roots have convergent small-parameter expansions to every fixed order. For finer spectral inverses, the horizontal error exponent tends to one from below, not to arbitrarily large values. Theorem~\ref{tdc:thm:marker} describes a local complex span marker without asserting a probability law.

\paragraph{Proof map.}
The source identification comes first. It justifies evaluating the span variable at one and determines the exact signed operator. Trace class gives meromorphic continuation. A finite exact determinant and a rational contour give a global zero count; separate real signs and a numerator test locate a genuine pole. A nonsingular resolvent differentiates the first-return quotient to find its residue. Finally, a boundary resolvent bound and Cauchy's estimate give \eqref{tdc:eq:error}. None of these steps can be replaced by observing a stable sequence of numerical roots.

\begin{remark}[{[write]} The OEIS entries, quoted, 7~October 2026]\label{tdc:rem:oeis}
The entries were read on 7~October 2026 in the OEIS internal format; the
quotations are verbatim.
\begin{enumerate}[label=(\alph*)]\raggedright
\item A394326~\cite{oeis} (revision \#15, 6~April 2026), offset $1$, author
 line ``\texttt{\_Morten Brydensholt\_, Mar 16 2026}'', is named
 \begin{center}\small
 \parbox{0.92\textwidth}{\raggedright\texttt{Total number of connected defect cluster types of weight n in the span-stratified cluster expansion for the row-word inversion statistic on standard Young tableaux of shape (m,m,m).}}
 \end{center}
 It prints forty terms, $1,2,1,1,3,5,\dots,25987091,42047032$, links a b-file
 ``\texttt{Table of n, a(n) for n = 1..100}'' (Sean A. Irvine) and
 Brydensholt's Python program, and has the two formula lines
 \begin{center}\small
 \parbox{0.92\textwidth}{\raggedright\texttt{a(n) \textasciitilde{} A * phi\^{}n where phi = (1+sqrt(5))/2, verified to 4 significant figures for n = 30..40.}\\[2pt]
 \texttt{a(n) \textasciitilde{} a(n-1) + a(n-2) for large n (empirical).}}
 \end{center}
 and the comment
 \begin{center}\small
 \parbox{0.92\textwidth}{\raggedright\texttt{The Fibonacci residual |a(n) - a(n-1) - a(n-2)|/a(n) < 0.0002 for n >= 30, confirmed by Hankel determinant analysis showing 2-state asymptotic dynamics.}}
 \end{center}
 Its first comment line states the hard-rod recursion of
 Section~\ref{tdc:sec:model} with $a(n)=\sum_{s\ge1}c(n,s)$.
\item A393486 (revision \#28) and A393487 (revision \#34), both by
 Brydensholt, 16~March 2026, are triangles of ``excess-tail polynomials'' of
 the same row-word inversion clusters; this report does not treat them.
\end{enumerate}
The b-file, fetched on 7~October 2026, agrees in all $100$ terms with the
write's own exact generator (Remark~\ref{tdc:rem:refuted}(d)). OEIS content is
published under CC~BY-SA~4.0. Nothing was submitted to the OEIS.
\end{remark}

\begin{remark}[{[write]} The two OEIS asymptotic formulas are false, 7~October 2026]\label{tdc:rem:refuted}
This records, under Vladimir's standing rule of 4~October 2026, the
source's refutation (Corollary~\ref{tdc:cor:empirical}) of the two formula
lines quoted in Remark~\ref{tdc:rem:oeis}(a), with its proof made explicit.
Write $\phi=(1+\sqrt5)/2$ and $\lambda=\rho^{-1}$.
\begin{enumerate}[label=(\alph*)]
\item \emph{``$a(n)\sim A\phi^n$'' is false for every constant $A$.} Since
 $a_n>0$ for $n\ge1$, the formula requires $A>0$ and $a_n/\phi^n\to A$. By
 Theorem~\ref{tdc:thm:main}, $a_n=c\rho^{-n}\{1+O((\rho/R)^n)\}$ with
 $c>0$, and $\rho>\rhol>0.618034>1/\phi$ (the second inequality is a
 comparison of decimals; the third is
 $0.618034^2+0.618034-1=6289/(25\cdot10^{10})>0$, Section~\ref{tdc:sec:inverse}).
 Hence $\rho\phi>1$ and $a_n/\phi^n=c(\rho\phi)^{-n}\{1+o(1)\}\to0$, so no
 $A>0$ works. Quantitatively $\lambda/\phi<1.6180189139/\phi<0.9999907$.
\item \emph{``$a(n)\sim a(n-1)+a(n-2)$'' is false.} Read as
 $(a_{n-1}+a_{n-2})/a_n\to1$ (equivalently, the residual is $o(a_n)$): by
 Theorem~\ref{tdc:thm:main}, $a_{n-j}/a_n\to\rho^j$, so the ratio tends to
 $\rho+\rho^2$, and by~\eqref{tdc:eq:empirical}
 $1.00001287<\rho+\rho^2<1.00001295$; it is not one.
\item \emph{What the observations were.} The OEIS observation is right as
 an observation and wrong as a law: $a_n/\phi^n$ decreases, but so slowly
 that $(\lambda/\phi)^{10}>0.99990$, invisible at four significant figures
 over $30\le n\le40$ (the write's exact terms give $a_n/\phi^n=0.18384\ldots$
 at $n=30$, $0.18375\ldots$ at $n=40$, $0.18329\ldots$ at $n=300$). The
 comment ``$|a(n)-a(n-1)-a(n-2)|/a(n)<0.0002$ for $n\ge30$'' is compatible
 with the theorem: the limit is $1-\rho-\rho^2=-0.0000129\ldots$; the write's
 exact terms give a largest value $0.000163\ldots$ (at $n=33$) over
 $30\le n\le300$; for all large $n$ it follows from the limit, but the
 range between is not certified here (Section~\ref{tdc:sec:open}). The
 ``Hankel determinant analysis'' is not assessed.
 \textup{(Added after the independent check of 7~October 2026: the range
 between is now closed. The check's exact terms $a_1,\dots,a_{1250}$ give
 $|a_n-a_{n-1}-a_{n-2}|/a_n<0.0002$ for $30\le n\le1250$ (largest value at
 $n=33$, as above), and the envelope~\eqref{tdc:eq:error} with the bounds
 \eqref{tdc:eq:rho}--\eqref{tdc:eq:c} gives it for every $n\ge982$, since
 with $t_n=(M/\cl)(\rhoh/R)^n$ the residual is at most
 $\{|1-\rho-\rho^2|+t_n(1+R+R^2)\}/(1-t_n)$, which is below $0.0002$ from
 $n=982$ on and decreases. So the comment's inequality holds for every
 $n\ge30$, resting on Theorem~\ref{tdc:thm:main} for $n\ge982$. Likewise
 $a_n/\phi^n$ decreases strictly for every $n\ge30$: exactly for
 $30\le n\le1250$, and from $n=1137$ on because
 $a_{n+1}/a_n\le\rho^{-1}(1+t_n\rho/R)/(1-t_n)<\phi$ there.)}
\item \emph{Trust boundary and corroboration.} (a) and (b) rest on the
 certified bracket~\eqref{tdc:eq:rho}, hence on the computer certificate of
 Sections~\ref{tdc:sec:global}--\ref{tdc:sec:local}, which the intake
 replayed (all certificates equal to the delivered ones). Independently of
 that certificate, the write computed $a_1,\dots,a_{300}$ exactly with its own
 code (the first-return walk of~\eqref{tdc:eq:edgesraw}, truncated by height;
 enlarging the truncation changes nothing) and $C=1-1/I$: the first forty
 equal the OEIS data and all $100$ b-file terms agree;
 $a_{300}/a_{299}=1.61801887065161\ldots$, inside~\eqref{tdc:eq:lambda}, and
 the residual at $n=300$ is $-1.29125514\ldots\cdot10^{-5}$, inside the band
 of~\eqref{tdc:eq:empirical}. This is corroboration, not proof.
\end{enumerate}
No claim of the source is refuted: the source itself states both failures.
\end{remark}

\subsection{Provenance, sources and reading conventions}\label{tdc:sec:provenance}
\begin{writenote}[Added 7 October 2026, batch~110]
\emph{Provenance.} This report is Report~183 of the session bundle that
ProveIt's \file{docs/incoming} intake processed as batch~110: the archive
\file{Tableau_Defect_Clusters_Certified_Asymptotics_Source.zip}
($2{,}104{,}583$ bytes, SHA-256 \file{442b90dea28f...95aef9676d04fb},
39 files, no wrapper directory), arrival commit \file{60f54ea06}, placed in
\file{8622ca7e5}. The text is the delivered \file{Report183.tex} (835 lines,
dated 3~October 2026; the delivered PDF has 21 pages), shipped as
\file{article.tex}. The programs are shipped under \file{code/} (the
delivered root \file{build.py}, \file{reproduce.py} and
\file{verify_manifest.py} among them), the certificates as
\file{data/certificates-*.json}, the generated summaries as
\file{data/generated-*.json}, the frozen data
(\file{N_H8.json}, \file{adjugate_H8.json}, \file{exact_coefficients.json} and
the $1{,}125{,}229$-byte \file{inverse_proposals.json}, an input of the
mandatory certificate) under \file{data/}, and \file{README_REPRODUCIBILITY.md}
at the report root (the report README maps every name). Not shipped, and
retrievable from the arrival commit: the delivered PDF, the checksum manifest
\file{SHA256SUMS.json} (38 entries, verified at placement), the delivered
README, replaced by the report README, and \file{certificates/coefficients.json},
a byte copy of \file{data/exact_coefficients.json}. The package records no
ProveIt commit; it cites one ProveIt report by its GitHub path and blob
(\cite{clustercomparison}). It carries no ``prepared for private review'' line
and no personal data. It is a single-source report, so no merge choice was
made. The write prefixed the 71 delivered labels with \file{tdc:} and
updated the 59 references to them; it added this subsection, labels on it and
on the subsection ``Useful directions left open''
(Section~\ref{tdc:sec:open}), the notes marked [write],
Remarks~\ref{tdc:rem:oeis}, \ref{tdc:rem:refuted}
and~\ref{tdc:rem:transseries}, and dated notes. The added remarks are the
last statements of their sections, this subsection follows the last delivered
text of Section~\ref{tdc:sec:main}, and the added displays are unnumbered, so
every theorem, section and equation of the source keeps its delivered number.

\emph{The sources, as the write read them} (7~October 2026).
\begin{itemize}\raggedright
\item OEIS A394326 with its b-file, and the names and authors of A393486 and
 A393487, in the internal format (Remark~\ref{tdc:rem:oeis}).
\item The transseries volume
 \path{Analysis/Transseries/docs/series-and-transseries/Transseries_And_Inversion/transseries_and_inversion.tex}:
 \file{p0:thm:lambert-core}, \file{p0:thm:core-reversion},
 \file{p0:def:three-inverses} and \file{p0:thm:staircase}, read for
 Remark~\ref{tdc:rem:transseries}.
\item The ProveIt report the source compares with,
 \path{SetTheory/Cardinals/docs/reports/log-concavity-and-unimodality/cluster-variable-log-concavity-counterexample/}:
 it concerns cluster variables of a cluster algebra, an unrelated use of the
 word ``cluster'', as the source says.
\item Not read by the write: the jOEIS program~\cite{joeis}, Brydensholt's
 Python program, Bornemann~\cite{bornemann,bornemannerratum},
 Simon~\cite{simon} and the ribbon-tableau manuscript~\cite{ribboncomparison};
 the source's account of them is reported, not confirmed.
\end{itemize}

\emph{What was checked at intake.} The batch-110 dossier read the
manuscript and found no error. The delivered \file{run_all.py} fails on
Windows only inside its two subprocess unit tests (CRLF text-mode output
compared with LF bytes); a driver running the same chain in-process
reproduced every certificate (global, local, remainder, inverse,
coefficients, independent global check, polynomial regeneration, local
tests) equal to the delivered files, in 22~s; \file{generate_proposals.py}
(mpmath~1.3.0) regenerated \file{inverse_proposals.json} in 19~s, identical
modulo CRLF. \emph{What the write checked} (Python~3.14.4, mpmath~1.3.0, own
code). Every proof was read line by line, with no error found; in exact
rational arithmetic it rechecked the bounds~\eqref{tdc:eq:lambda} from the
endpoints of~\eqref{tdc:eq:rho}, the value $6289/(25\cdot10^{10})$, the band
of~\eqref{tdc:eq:empirical} at both endpoints, $Q_{600}<0.138$ and
$Q_{600}(1+R^{-1})<0.356406<\rhoh^{-1}-1$ of~\eqref{tdc:eq:growthgate}, and
all $125$ digits of $Y_0$~\eqref{tdc:eq:Y0}. Its own exact generator
(Remark~\ref{tdc:rem:refuted}(d)) gives $\rho=0.618039763403548969\ldots$ and
$c\approx0.183811179648135$, agreeing with the source's diagnostic
refinements, and shows $a_{n+1}>a_n$ for $4\le n<300$. The write did not rerun
the certificate. (Added after the independent check of 7~October 2026: this
sentence disagrees with the report README and the write's commit message,
which record a Route~B replay at the write, all eight certificates and
\file{N_H8.json} equal, in 45~s. Either way the chain has been replayed
independently of the write: at placement and again by the check, with the
same result; see the note before the bibliography.)

\emph{Relation to the repository.} Before batch~110 no file of the
repository named A394326, A393486 or A393487, and no report treated
tableau inversion clusters. No statement of this report is formalized in
Lean or Rocq, and its place in the collection confers no formal status. The
cluster-variable report cited by the source shares only the word
``cluster''; no neighbouring report shares a result, so no reciprocal note is
proposed. The inverses are compared with the transseries volume in
Remark~\ref{tdc:rem:transseries}.
\end{writenote}

\begin{writenote}[What is not claimed, collected from the source]
The sequence, its extraction and the hard-rod recursion are Brydensholt's;
ballot paths, renewal, Fredholm determinants, coefficient extraction and
Lagrange inversion are classical; the bounded source screen ``is not a
worldwide priority determination''. The identification of the model does not
prove that every cluster coefficient is nonnegative. No infinite diagonal
gauge is claimed bounded, and no positive-operator (Perron--Frobenius)
argument is used for the signed operator. The numerical refinements of
$\rho$ and $c$ beyond~\eqref{tdc:eq:rho}--\eqref{tdc:eq:c} are not
certified; $M=2500$ is a transparent rather than sharp envelope. Only the
dominant pole is certified: Theorem~\ref{tdc:thm:spectral} gives no
convergent pole sum, no uniformity as $r\uparrow1$ and no simplicity of later
zeros. The Lagrange series expand bounding roots, not the step function; the
spectral inverse constants are existential and do not give errors
$O(y^{-K})$, $K>1$. Theorem~\ref{tdc:thm:marker} certifies no marker radius,
derivative or variance, and implies no probability law. The certificates are
not a machine-checked proof; hashes are not mathematical evidence. The write
adds: its checks are exact finite or floating computations; it claims no
novelty for any inversion.
\end{writenote}

\medskip
\begingroup\small
\noindent\emph{Reading conventions} (letters the source reuses, with the
tempting false readings; no symbol of the source was renamed; symbols of the
transseries volume carry the subscript ``vol'' in
Remark~\ref{tdc:rem:transseries}).\par\smallskip
\renewcommand{\arraystretch}{1.1}
\begin{longtable}{@{}>{\raggedright\arraybackslash}p{0.95in}>{\raggedright\arraybackslash}p{3.3in}>{\raggedright\arraybackslash}p{1.85in}@{}}
\toprule
Symbol & Meaning here (where) & Not to be read as\\
\midrule
\endhead
$C$ & $C(q)=\sum a_nq^n$; the bivariate cluster series $C(t,q)$ (Sections~\ref{tdc:sec:model}, \ref{tdc:sec:marker}); $C_0$ a block of~\eqref{tdc:eq:blocks} & each other\\
$c$, $c_{w,s}$, $c(t)$ & the amplitude~\eqref{tdc:eq:c}; cluster coefficients; the marked amplitude & each other\\
$K$, $L$, $L$ & the three-step operator; the corrected operator~\eqref{tdc:eq:Landdet}; but $L=\log\lambda$ in Section~\ref{tdc:sec:inverse} & each other; the volume's $L_{\mathrm{vol}}$\\
$D$, $N$, $N_8$, $n$ & the determinants~\eqref{tdc:eq:Landdet} (also $D(t,q)$, $N(t,q)$); the degree-161 core polynomial; an index & each other\\
$R$, $\Res$, $r$, $r_*$ & the radius $63/100$; a residue; a fixed spectral radius, or $|z|$; $4/5$ & each other\\
$M$, $M_r$ & $2500$; the constant of Theorem~\ref{tdc:thm:spectral} & each other\\
$A$, $B$, $E$ & states $(1,0)$, $(0,1)$, $(0,2)$; but $A_0$, $B_0$, $C_0$ blocks, $A_8$, $B_8$ core matrices, $B$ an inverse proposal, $A_n(t)$ marked coefficients, $A_r(t)$ a spectral carrier, $E_{r,n}$ a remainder & each other; the volume's $A_n$, $\mathcal A_{\mathrm{vol}}$\\
$I$, $\Id$, $J$, $F$, $F_m$, $Z$ & the first-return series $J(1,q)$; the identity; first returns; the deficit series; the deficit polynomial; the source's $Z_n$ & each other\\
$G$, $g$, $g_0$ & the potential $G(a,b,c)$, but $G=C-c/(1-q/\rho)$ in Section~\ref{tdc:sec:error}; the gauge $g(x,y)$; the vector~\eqref{tdc:eq:Cresolvent} & each other\\
$h$, $H_\zeta$, $H_r$ & the vector $e_A+z^3e_E$, but the height $x+y$ in Lemma~\ref{tdc:lem:trace}; the local factors and remainder of Theorem~\ref{tdc:thm:spectral} & each other\\
$\kappa$, $\delta$, $\tau$, $\eta$ & the inverse bound and Schur constants of Section~\ref{tdc:sec:global}; $\eta=\kappa\delta$ & the volume's $\kappa_{\mathrm{vol}}$\\
$\lambda$, $\rho$, $\phi$ & $\rho^{-1}$; the pole; the golden ratio & the volume's $\rho_{\mathrm{vol}}$\\
$\alpha$, $\beta$, $\mu$, $x_0$, $x_\pm$ & the inverse parameters~\eqref{tdc:eq:inverseparams}; the bounding roots & the volume's $\mu_{\mathrm{vol}}$\\
$x$, $y$ & gap coordinates (also $x=0.618034$); but $x$ real and $y$ a threshold in Section~\ref{tdc:sec:inverse} & each other\\
$u$, $v$, $w$ & nodes $u_j$, and $u(z)$; $v=z^3e_E$, but $v=e^{L(x-x_0)}$ in the proof of Proposition~\ref{tdc:prop:lagrange}; a weight, and $w=v-1$ & each other\\
$t$, $\epsilon$, $\varepsilon$ & the span marker; $\pm1$ in the proof of Proposition~\ref{tdc:prop:lagrange}, a marker radius in Theorem~\ref{tdc:thm:marker}; a residual & the volume's $t$\\
$S$, $\Sigma$, $\mathcal T_s$, $T_0$ & the state set; the Schur complement; the homotopy operator; a tail block & each other\\
$P_j$, $\PP_r$, $Q_n$ & $\binom{t+j-1}{j-1}$; the set of poles; $(M/\cl)(\rhoh/R)^n$ & each other\\
\bottomrule
\end{longtable}
\endgroup


\section{From the source extraction to a first-return quotient}
\label{tdc:sec:model}
\subsection{The inversion deficit as a nonnegative path weight}
A standard Young tableau of shape $(m,m,m)$ is encoded by reading its entries $1,\ldots,3m$ and recording their row labels. Its row word contains $m$ copies of each of $1,2,3$; every prefix has counts $a\ge b\ge c$. Conversely, each such ballot word gives exactly one standard tableau by placing the labels successively in their indicated rows.

The inversion statistic in the source increases by $b+c$ on appending a $1$, by $c$ on appending a $2$, and by zero on appending a $3$. Introduce
\[
 G(a,b,c)=ab+ac+bc-\frac{b(b+1)}2-c(c+1),\qquad
 x=a-b,\quad y=b-c.
\]
Subtracting the inversion increment from the increment of $G$ gives the following weighted gap walk in $\NN^2$:
\begin{equation}\label{tdc:eq:edgesraw}
\begin{array}{c|c|c}
 \text{letter}&\text{new gap state}&\text{deficit cost}\\ \hline
 1&(x+1,y)&0\\
 2&(x-1,y+1),\quad x>0&x-1\\
 3&(x,y-1),\quad y>0&x+2y-2
\end{array}
\end{equation}
All allowed costs are nonnegative. At $(a,b,c)=(m,m,m)$ the potential is $3\binom m2$. The word $(123)^m$ has cost zero. Hence the maximal inversion degree is exactly $3\binom m2$, and the source's reversed top coefficients count these walks by total deficit. Let $F_m(q)$ be that deficit polynomial, with $F_0(q)=1$.

\subsection{First returns and a finite span at every weight}
Let
\[
 F(t,q)=\sum_{m\ge0}F_m(q)t^m,\qquad
 J(t,q)=\sum_{\text{primitive returns}}t^m q^W.
\]
A primitive return leaves $(0,0)$ and does not revisit it before its final step; $m$ is its common final letter count and $W$ its deficit. Concatenating balanced blocks adds deficit. This follows either from \eqref{tdc:eq:edgesraw} or from the fact that two blocks of sizes $m_1,m_2$ create $3m_1m_2$ cross inversions, exactly the cross term of $3\binom{m_1+m_2}2$. The unique decomposition at successive returns therefore gives
\begin{equation}\label{tdc:eq:renewal}
 F=\frac1{1-J}.
\end{equation}

\begin{lemma}[The span bound]\label{tdc:lem:span}
A primitive return of size $m$ has deficit $W\ge m-1$. Thus
\[
 J(t,q)/t=1+\sum_{w\ge1}\sum_{1\le s\le w}j_{w,s}q^w t^s.
\]
The same support restriction $1\le s\le w$ holds for $1-t/J(t,q)$.
\end{lemma}
\begin{proof}
There are $m$ letters $3$. A $3$ step has cost zero only from $(0,1)$, where it would return to the origin. Each of the first $m-1$ occurrences must therefore have cost at least one. The unique size-one path is $123$, of deficit zero. Dividing by $t$ replaces $m$ by $s=m-1$, so all other monomials have $1\le s\le w$. Products preserve this cone of exponents, and the formal geometric-series expansion of the reciprocal proves the final assertion. The same reasoning shows the only zero-deficit complete word is $(123)^m$.
\end{proof}

\subsection{Why the finite source extraction gives this exact quotient}
The source uses $Z_n(q)=F_{n+1}(q)$ and a cluster series
\[
 C(t,q)=\sum_{w,s\ge1}c_{w,s}q^w t^s.
\]
Its hard-rod recursion is $Z_0(q)=1$ and
\[
 Z_n(q)=Z_{n-1}(q)+\sum_{\substack{w\ge1\\1\le s\le n}}
 c_{w,s}q^w Z_{n-s}(q).
\]
Consequently, with $Z(t,q)=\sum_{n\ge0}Z_n(q)t^n$,
\[
 Z=\frac{F-1}{t}=\frac{J}{t(1-J)},\qquad
 Z=\frac1{1-t-C}.
\]
Eliminating $Z$ gives the exact identity
\begin{equation}\label{tdc:eq:sourcequotient}
 C(t,q)=1-\frac{t}{J(t,q)}.
\end{equation}
This also explains the source's sequential finite extraction. Suppose the cluster coefficients of weights less than $w$ have been inserted into the recursion, and denote the resulting series by $Z_{<w}$. Since $C$ has no weight-zero term,
\[
 [q^w](Z-Z_{<w})=\frac{C_w(t)}{(1-t)^2},\qquad
 C_w(t)=[q^w]C(t,q).
\]
If $r_n=[t^nq^w](Z-Z_{<w})$, two backward differences recover
$c_{w,s}=r_s-2r_{s-1}+r_{s-2}$, with negative subscripts interpreted as zero. This is precisely the inspected implementation's residual, first-difference, and cluster update. Lemma~\ref{tdc:lem:span} proves $c_{w,s}=0$ for $s>w$, so stopping its span sum at $s=w$ loses nothing. The code's integer storage details are an implementation, not a restriction on this mathematical definition.

All these identities first hold formally in the two variables. The support lemma makes every coefficient of $q^w$ in $J/t$ and $C$ a polynomial in $t$. Thus specialization at $t=1$ is legitimate coefficient by coefficient. Write $I(q)=J(1,q)$. Then
\begin{equation}\label{tdc:eq:CfromI}
 C(q)=1-\frac1{I(q)},\qquad I(0)=1.
\end{equation}
This establishes the model for all indices, beyond any finite agreement test. It does not prove that every extracted cluster coefficient is nonnegative. Although the initial first-return coefficients are nonnegative, a zero of its meromorphic continuation need not lie in the disk of convergence of that positive series.

\section{The signed trace-class determinant}
\label{tdc:sec:operator}
\subsection{A finite-path gauge and period-three compression}
Delete the origin from the gap graph. Set $q=z^2$ and
\[
 g(x,y)=x^2+xy+y^2-2x-2y.
\]
On each remaining edge $u\to v$, multiply $q^d$ by $z^{g(v)-g(u)}$. Its exponent becomes, for letters $1,2,3$ respectively,
\begin{equation}\label{tdc:eq:gauged}
 2x+y-1,\qquad x+y-1,\qquad x+2y-1.
\end{equation}
These exponents are nonnegative on every allowed edge off the origin. The gauge factors telescope along a finite path. We use this path identity; we do not claim that an infinite diagonal gauge is bounded or that it gives a bounded similarity of infinite operators.

Each step advances $x+2y$ by one modulo three. Define
\[
 S=\{(x,y)\in\NN^2:x+2y\equiv1\pmod3\},\quad
 A=(1,0),\quad B=(0,1),\quad E=(0,2).
\]
Let $K(z)$ be the three-step transition matrix on $S$ in the origin-deleted graph. Rows are starting states and columns are ending states. Every entry sums the monomial weights of all three-step paths, including multiplicities.

A primitive word begins with $1$ and ends with $3$. Its outer two gauge exponents are $-1$ and $1$, so cancel. The intervening path runs from $A$ to $B$. Its last step comes either from $A$ with weight one, or from $E$ with weight $z^3$. A primitive size-$m$ path consists of $m-1$ compressed moves followed by that last step. Therefore, near zero,
\begin{equation}\label{tdc:eq:resolventI}
 I(z^2)=e_A^{\mathsf T}(\Id-K(z))^{-1}h(z),\qquad
 h(z)=e_A+z^3e_E.
\end{equation}
The convergence justification follows below; formally the identity is already the first-return enumeration.

\subsection{Trace class and removal of the square-root branch}
\begin{lemma}\label{tdc:lem:trace}
For $|z|<1$, $K(z)$ is trace class on $\ell^2(S)$ and is holomorphic in trace norm. Its principal height sections converge locally uniformly in trace norm. Moreover $K(0)=0$.
\end{lemma}
\begin{proof}
Put $h=x+y$. Every one-step exponent in \eqref{tdc:eq:gauged} is at least $h-1$, and height falls by at most one per step. Hence each three-step exponent from height $h$ is at least $3h-6$. There are at most $27$ such paths. For each fixed $0<r<1$, the entrywise absolute sum over states of height at least two is bounded by
\begin{equation}\label{tdc:eq:nuclear}
 \sum_{h\ge2}27(h+1)r^{3h-6}<\infty.
\end{equation}
There are at most $h+1$ states at height $h$, and the height-one contribution is finite. Expanding in rank-one matrix units, each of trace norm one, proves trace class and locally uniform trace-norm convergence. The polynomial matrix-unit series is normally convergent, giving trace-norm holomorphy. Equivalently, Cauchy estimates on a slightly larger subdisk control its derivatives. Removing all entries outside a principal height section has trace norm tending to zero by the same absolute sum.

For height at least three the bound on exponents is positive. The only lower states in $S$ are $A,E$; direct enumeration shows their three-step exponents are all positive. Thus $K(0)=0$. In particular its norm is less than one sufficiently near zero, which justifies the analytic Neumann series in \eqref{tdc:eq:resolventI}.
\end{proof}

Set $\OO=S\setminus\{A\}$ and $v=z^3e_E$ on this space, and use $A/\OO$ blocks. Define
\begin{equation}\label{tdc:eq:Landdet}
 L(z)=K_{\OO\OO}(z)-vK_{A\OO}(z),\qquad
 D(z^2)=\det(\Id-K(z)),\quad N(z^2)=\det(\Id-L(z)).
\end{equation}
The row $K_{A\OO}$ has finite support, so the correction is rank one and polynomial. Thus $L$ is also trace class and trace-norm holomorphic, with $L(0)=0$.

For each gauged path the parity of its exponent is $g(v)-g(u)$. The diagonal operator $U_{u,u}=(-1)^{g(u)}$ is bounded with bounded inverse, and
\[
 K(-z)=UK(z)U^{-1},\qquad L(-z)=U_{\OO}L(z)U_{\OO}^{-1}.
\]
The correction has the same parity since $g(A)=-1$, $g(E)=0$. Determinant invariance under bounded similarity makes both determinants even and holomorphic in $z$. Their even power series define holomorphic functions $D,N$ of $q=z^2$ on $|q|<1$, independently of any square-root branch; $D(0)=N(0)=1$.

\begin{proposition}[Exact determinant quotient]\label{tdc:prop:quotient}
The meromorphic continuation of the source cluster series is
\begin{equation}\label{tdc:eq:detquotient}
 I(q)=\frac{N(q)}{D(q)},\qquad C(q)=1-\frac{D(q)}{N(q)}\qquad(|q|<1).
\end{equation}
All rows of $L$ except the $E$ row are those of $K_{\OO\OO}$. The exceptional nonzero entries are
\begin{equation}\label{tdc:eq:exceptional}
 L_{E,E}=z^8,\quad L_{E,(3,2)}=z^9,\quad
 L_{E,(1,3)}=z^5+z^7,\quad L_{E,(4,0)}=-z^{12}.
\end{equation}
\end{proposition}
\begin{proof}
In a finite principal section, replace column $A$ of $\Id-K$ by $h=e_A+v$. The resulting matrix is
\[
 \begin{pmatrix}1&-K_{A\OO}\\v&\Id-K_{\OO\OO}\end{pmatrix}.
\]
Subtract $v$ times its first row from the other rows. Its determinant becomes $\det(\Id-K_{\OO\OO}+vK_{A\OO})=\det(\Id-L)$. Cramer's rule proves the first identity in a neighborhood of zero. The finite operators and determinants converge in trace norm and locally uniformly, while resolvents converge near zero by their Neumann series. Passing to the limit proves the identity there. Equation~\eqref{tdc:eq:CfromI} proves the second one, and the identity theorem gives the meromorphic continuation. Finally, enumerating the at most $27$ three-step paths from $E$ and subtracting $z^3$ times the $A$ row gives \eqref{tdc:eq:exceptional}.
\end{proof}

The negative entry in \eqref{tdc:eq:exceptional} is essential. A positive-operator Perron--Frobenius argument for $L$ would not be justified. The poles of $I=N/D$ can also disappear when forming $C=1-D/N$; its relevant singularities come from zeros of $N$ that survive cancellation with $D$.

\paragraph{Classical analytic facts used.}
For trace-class operators, the determinant is continuous under trace-norm approximation, is holomorphic for a trace-norm holomorphic family, is invariant under bounded similarity, and is nonzero exactly when the identity minus the operator is invertible. Its multiplicativity supplies the Schur factorization below. These are the general determinant facts in Bornemann~\cite[Sections 2--4]{bornemann}, especially equations (3.1)--(3.5) and the general perturbation inequality (4.1); see also Simon~\cite{simon}. The target's trace-class hypothesis has been proved directly in Lemma~\ref{tdc:lem:trace}. No positive self-adjoint estimate is used for this signed operator.

\section{An exact global zero count}
\label{tdc:sec:global}
\subsection{Schur bounds and the two normed spaces}
Take a finite height core and split either $\Id-K$ or $\Id-L$ as
\begin{equation}\label{tdc:eq:blocks}
 \begin{pmatrix}A_0&-B_0\\-C_0&\Id-T_0\end{pmatrix}.
\end{equation}
Here $A_0$ denotes the finite identity-minus-core matrix. Suppose the induced $\ell^\infty$ row norms obey
\[
 \norm{B_0}_\infty\le b,\quad \norm{C_0}_\infty\le d,
 \quad\norm{T_0}_\infty\le\tau<1.
\]
These are not estimates of the $\ell^2$ operator norm. Nevertheless they give the required determinant calculation on $\ell^2$. A vector in $\ell^2$ is bounded, so a nonzero $\ell^2$ kernel vector of $\Id-T_0$ would contradict its $\ell^\infty$ Neumann inverse. Since $T_0$ is compact on $\ell^2$, the Fredholm alternative gives invertibility on $\ell^2$. For a square-summable right side, its inverse agrees with the unique bounded Neumann-series solution. Therefore
\begin{equation}\label{tdc:eq:schur}
 \Sigma=B_0(\Id-T_0)^{-1}C_0,
 \qquad\norm{\Sigma}_\infty\le\delta:=\frac{bd}{1-\tau}
\end{equation}
is valid for the finite Schur complement.

If $\norm{A_0^{-1}}_\infty\le\kappa$ and $\kappa\delta<1$, then $A_0-s\Sigma$ is invertible for $0\le s\le1$. Scaling the upper-right block in \eqref{tdc:eq:blocks} by $s$ gives
\begin{equation}\label{tdc:eq:fredholmschur}
 \det(\Id-\mathcal T_s)=\det(\Id-T_0)\det(A_0-s\Sigma),
\end{equation}
where $\mathcal T_s$ denotes the entire trace-class block operator being subtracted from the identity, not just its tail block. The factorization follows by block triangular elimination and determinant multiplicativity; the triangular factors differ from the identity by finite-rank off-diagonal operators and have determinant one. On a complex boundary this homotopy preserves the zero count. At a real point it preserves the sign. The real tail determinant is positive because scaling the tail itself from zero to $T_0$ preserves nonvanishing and starts at one.

\subsection{The finite determinant is genuinely a determinant}
Use $\OO_8=\{u\in\OO:x+y\le8\}$, ordered by increasing height and then increasing $x$. Its complete ordered state list is
\begin{align*}
 &(0,2),(2,1),(1,3),(4,0),(0,5),(3,2),(2,4),\\
 &(5,1),(1,6),(4,3),(7,0),(0,8),(3,5),(6,2).
\end{align*}
Let $A_8(z)=\Id-L_8(z)$. The supplied integer polynomial $N_8(q)$ has degree $161$. The exact verifier checks the complete state list, reconstructs every matrix entry from the transition rules, and checks the full polynomial identity
\begin{equation}\label{tdc:eq:adjugate}
 A_8(z)\operatorname{Adj}_8(z)=N_8(z^2)\Id_{14}.
\end{equation}
This identity is not used alone to identify the determinant: an extraneous common polynomial factor could multiply both the proposed adjugate and denominator.

For an independent determinant check, construct the ungauged $q$-matrix from \eqref{tdc:eq:edgesraw}, compress three steps, and subtract the correction $q^2e_E K^{\mathrm{raw}}_{A\OO}$. Write its identity-minus-core matrix as $B_8(q)$. The correction exponent is correct because
\[
 2\cdot2+g(A)-g(E)=3.
\]
Entrywise finite gauge conjugacy, checked including the signed row, gives
$\det A_8(z)=\det B_8(z^2)$ for $z\ne0$, hence identically as polynomials. The row maximum degrees of $B_8$ are
\[
 (4,6,10,9,18,12,18,14,27,19,18,36,27,21),
\]
whose sum is $239$. Thus $\deg\det B_8\le239$ by the Leibniz expansion. At each of the $240$ distinct integers $q=0,1,\ldots,239$, exact fraction-free Bareiss elimination gives
\[
 \det B_8(q)=N_8(q).
\]
Each division is explicitly checked to have zero remainder; row swaps track their signs and singular cases are handled. The difference has degree at most $239$ and $240$ roots, so is identically zero. This is a degree-bounded identity proof, not extrapolation from a finite fit. It proves $N_8(z^2)=\det A_8(z)$ without a symbolic algebra dependency.

\subsection{Rational parametrization of the whole circle}
Set $m_*=2048$ and, for $0\le j<m_*$, define
\begin{equation}\label{tdc:eq:rationalmesh}
 u_j=\frac{2j+1}{2m_*},\qquad
 q_j=R\frac{1-u_j^2+2iu_j}{1+u_j^2}.
\end{equation}
Follow these points by their rotations through $\pi/2$, $\pi$, and $3\pi/2$, retaining the same within-quadrant order and closing the last edge to the first. This gives $8192$ rational complex points in positive circular order. With $s=2j+1$, $t=2m_*$, the unrotated point is $(a+ib)/d$, where
\[
 a=63(t^2-s^2),\quad b=126st,\quad d=100(t^2+s^2)>0.
\]
The verifier checks $(a^2+b^2)100^2=d^2 63^2$ exactly at every node.

For $w(u)=(1-u^2+2iu)/(1+u^2)$ one has $|w(u)|=1$ and
$w'(u)/w(u)=2i/(1+u^2)$. Thus the angle strictly increases with speed at most two. Interior parameter gaps have length $1/m_*$. Across a quadrant join the omitted pieces have total length $1/m_*$ as well; this includes the final-to-first join. Every corresponding circle arc therefore has length at most $2R/m_*$.

For a node $(a+ib)/d$, integer Horner evaluation represents a polynomial partial value as $(X+iY)/D$ with positive denominator. Multiplication by the node and addition of coefficient $p$ updates it by
\[
 (X,Y,D)\longmapsto(Xa-Yb+pDd,\ Xb+Ya,\ Dd).
\]
The final denominator is $d^{161}$. At each node the certificate verifies
\begin{equation}\label{tdc:eq:nodemod}
 625(X_j^2+Y_j^2)>D_j^2,\qquad Y_j\ne0,
\end{equation}
so $|N_8(q_j)|>1/25$ and no image vertex lies on the real axis.

Writing $N_8(q)=\sum p_kq^k$, the exact coefficient majorant gives
\[
 \sum_{k=1}^{161}k|p_k|R^{k-1}<24.
\]
The image of any successive circle arc stays within distance
\[
 24\frac{2R}{m_*}=\frac{189}{12800}
\]
of its first image vertex. Therefore, on the entire circle,
\begin{equation}\label{tdc:eq:wholecircle}
 |N_8(q)|>\frac1{25}-\frac{189}{12800}
 =\frac{323}{12800}>\frac2{125}.
\end{equation}
Both the image arc and its straight chord lie in the same convex disk avoiding zero. Linear interpolation in these disks supplies an endpoint-fixing zero-avoiding homotopy from the image curve to its polygon. Their winding numbers agree.

\subsection{The polygon winds once}
For adjacent image values $(X_1+iY_1)/D_1$ and $(X_2+iY_2)/D_2$, opposite signs of $Y_1,Y_2$ characterize a real-axis crossing. Its intercept is
\begin{equation}\label{tdc:eq:crossing}
 x_{\mathrm{cross}}=
 \frac{X_1Y_2-X_2Y_1}{Y_2D_1-Y_1D_2}.
\end{equation}
This formula retains the potentially unequal denominators. The verifier excludes a zero numerator or denominator. An upward positive-ray crossing contributes $+1$ and a downward one $-1$. The complete exact crossing list, with zero-based edge indices, is
\begin{center}
\begin{tabular}{rll@{\qquad}rll}
\toprule
Edge&Direction&Ray&Edge&Direction&Ray\\
\midrule
673&up&positive&2092&down&positive\\
3193&up&positive&4095&down&positive\\
4673&up&positive&6056&down&positive\\
7177&up&positive&8191&down&negative\\
\bottomrule
\end{tabular}
\end{center}
The positive-ray total is $4-3=1$. The closing edge is included. By the argument principle, $N_8$ has exactly one zero in $|q|<R$, counted with multiplicity.

\subsection{The infinite tail cannot change the count}
All following coefficient majorants use $|z|\le r_*=4/5$, which contains the lifted boundary since $R<r_*^2$. Enumerate the exact signed polynomial rows through height $22$. A three-step transition changes height by at most three; hence no tail-to-core row beyond height $11$ contributes. For all unenumerated rows, of height $h\ge23$, the decreasing bound $27r_*^{3h-6}$ applies. The exceptional row is in the core and was checked separately. The exact maxima satisfy
\begin{align}
 b&=\frac{194073697907113984}{7450580596923828125}<\frac{27}{1000},\notag\\
 d&=\frac{147844731517272064}{7450580596923828125}<\frac{21}{1000},\label{tdc:eq:tailnumerics}\\
 \tau&=\frac{34489026312519942504710144}{1818989403545856475830078125}
 <\frac1{50}.\notag
\end{align}
The adjugate coefficient majorant is $\norm{\operatorname{Adj}_8(z)}_\infty<63/10$. Equations~\eqref{tdc:eq:adjugate} and \eqref{tdc:eq:wholecircle} give the conservative bound
\begin{equation}\label{tdc:eq:kappadelta}
 \kappa=\frac{1575}{4},\qquad \delta=\frac{81}{140000},
 \qquad \kappa\delta=\frac{729}{3200}<1.
\end{equation}
The stronger circle lower bound also permits $\kappa=80640/323$, but retaining \eqref{tdc:eq:kappadelta} simplifies the subsequent certified estimates.

Apply the Schur homotopy on $|z|=\sqrt R$. At $s=0$ its determinant is $N_8(z^2)$ times a zero-free tail determinant; it therefore has two $z$ zeros. No member of the homotopy has a boundary zero. Parity persists when a cross block is scaled, and the determinant has constant term one. Thus the full determinant has two $z$ zeros, equivalently one $q$ zero, counted with multiplicity. Real coefficients make the latter real and multiplicity one makes it simple. This proves the global zero-count part of Theorem~\ref{tdc:thm:main}; positive location and noncancellation still require the next section.

\section{Real location and the residue}
\label{tdc:sec:local}
\subsection{What the local interval arithmetic certifies}
Use height $20$, with $76$ states for $L$ and $77$ for $K$, and the exact real endpoints
\begin{equation}\label{tdc:eq:zinterval}
 z_-=0.78615504,\qquad z_+=0.78615506.
\end{equation}
Their squares are precisely $\rhol,\rhoh$ in \eqref{tdc:eq:rho}. Each local interval has endpoints that are integers divided by $2^{220}$. Addition is exact; products and quotients are rounded outward by integer floor and ceiling. Division rejects intervals containing zero. Exponentiation uses interval multiplication.

For a finite interval matrix $A_0$, an exact dyadic proposal $B$ for its inverse is accepted only after the outward residual proves
\[
 \varepsilon=\norm{\Id-BA_0}_\infty<1.
\]
Then the finite-dimensional Neumann argument proves
\begin{equation}\label{tdc:eq:inversecertificate}
 \norm{A_0^{-1}}_\infty\le\frac{\norm B_\infty}{1-\varepsilon}.
\end{equation}
The proposal is not assumed to enclose the inverse. Interval Gaussian elimination separately certifies determinant signs, rejecting any pivot that is not separated from zero. Proposed pivot choices affect efficiency only, never acceptance. The public package freezes exact dyadic proposals, so replay uses only standard-library integer arithmetic; optional high-precision numerical generation of proposals is outside the proof's trust requirement.

The tail estimates again enumerate finite-band interactions and bound all later rows by $27r^{3h-6}$. With $r=z_+$, the explicit Schur tests give
\begin{center}
\begin{tabular}{lll}
\toprule
Finite determinant&Certified conclusion&Schur product bound\\
\midrule
$N_{20}(z_-^2)$&$>6.26786\cdot10^{-8}$&$<0.000167$\\
$N_{20}(z_+^2)$&$<-5.68766\cdot10^{-8}$&$<0.000184$\\
$D_{20}(z^2)$, $z\in[z_-,z_+]$&$\in[-0.431905973,-0.431905414]$&$<3.78\cdot10^{-11}$\\
\bottomrule
\end{tabular}
\end{center}
These displayed determinant intervals concern the finite sections. The Schur sign homotopy transfers their signs, not these same numerical intervals, to the full determinants. Thus $N(z_-^2)>0$, $N(z_+^2)<0$, and $D(z^2)<0$ throughout the real bracket. The intermediate value theorem locates the real zero. The global count proves it is the unique simple zero $\rho$, and the test for $D$ excludes cancellation. Since $N$ crosses from positive to negative at its sole simple zero, $N'(\rho)<0$.

\subsection{Differentiation through a nonsingular resolvent}
Near the bracket, $\Id-K$ stays invertible. Define
\[
 u(z)=(\Id-K(z))^{-1}h(z),\qquad I(z^2)=u_A(z).
\]
Differentiate the exact operator equation:
\begin{equation}\label{tdc:eq:du}
 u'(z)=(\Id-K(z))^{-1}\{K'(z)u(z)+h'(z)\}.
\end{equation}
This avoids differentiating a nearly singular determinant or any asymptotic remainder.

Here are the explicit tail estimates used to certify \eqref{tdc:eq:du}. Let $b_1,d_1,\tau_1$ bound the derivatives of the three blocks in \eqref{tdc:eq:blocks}. Differentiating the Schur correction gives
\begin{equation}\label{tdc:eq:deltaslope}
 \norm{\Sigma'}_\infty\le\delta_1:=%
 \frac{b_1d}{1-\tau}+\frac{b\tau_1d}{(1-\tau)^2}
 +\frac{bd_1}{1-\tau}.
\end{equation}
These derivative majorants sum $|p|e r^{e-1}$ for each monomial $pz^e$. Beyond the enumerated rows, $e\ge3h-6$. The inequality $(e+1)r<e$ guarantees that $e r^{e-1}$ decreases; this inequality is explicitly checked at the first omitted exponent before using the bound for all later exponents.

Let $u_f,u_f'$ denote the exact finite-core solutions, enclosed by residual-certified interval solves. Put $\eta=\kappa\delta<1$ and let $u_c,u_c'$ denote the full solution restricted to the core. The tail right side is zero because $h$ is supported on $A,E$. The finite and full Schur equations imply
\begin{align}
 \norm{u_c-u_f}_\infty&\le
 \frac{\kappa\delta}{1-\eta}\norm{u_f}_\infty=:e_0,\label{tdc:eq:uerror}\\
 \norm{u_c'-u_f'}_\infty&\le\frac{\kappa}{1-\eta}
 \bigl(\delta\norm{u_f'}_\infty+
 \norm{K_{\mathrm{core}}'}_\infty e_0+
 \delta_1(\norm{u_f}_\infty+e_0)\bigr).
 \label{tdc:eq:duerror}
\end{align}
For completeness, subtracting $(A_0-\Sigma)u_c=h$ from $A_0u_f=h$ gives the first estimate. Differentiating and subtracting the two equations yields
\[
 (A_0-\Sigma)(u_c'-u_f')=
 \Sigma u_f'+K_{\mathrm{core}}'(u_c-u_f)+\Sigma'u_c,
\]
which proves the second one. The certified derivative tail error is below $8.259\cdot10^{-9}$.

Combining these bounds with the finite solve gives, at $z=\sqrt\rho$,
\begin{align}
 13.8403718617&<\frac{d}{dz}I(z^2)<13.8405060269,\label{tdc:eq:Iz}\\
 8.8025712521&<I'(\rho)<8.8026568060.\label{tdc:eq:Iq}
\end{align}
Since $I(\rho)=0$ and $C=1-1/I$, its residue and normalized amplitude are
\begin{equation}\label{tdc:eq:amplitude}
 \Res_{q=\rho}C=-\frac1{I'(\rho)},\qquad
 c=\frac1{\rho I'(\rho)}
 =\frac{2}{\sqrt\rho\,\frac{d}{dz}I(z^2)|_{z=\sqrt\rho}}>0.
\end{equation}
Outward evaluation proves \eqref{tdc:eq:c}. Differentiating $I=N/D$ at its zero also gives $I'(\rho)=N'(\rho)/D(\rho)$, agreeing with \eqref{tdc:eq:exactconstants}.

\section{A uniform and explicit coefficient error}
\label{tdc:sec:error}
A second block identity is particularly effective for estimating $C$ on the global circle:
\begin{equation}\label{tdc:eq:Cresolvent}
 C=K_{AA}+K_{A\OO}(\Id-L)^{-1}g_0,
 \qquad g_0=K_{\OO A}+z^3e_E(1-K_{AA}).
\end{equation}
To derive it near zero, divide the $\OO$ block of the equation
$(\Id-K)u=h$ by $u_A=I$. If $x=u_{\OO}/u_A$, then
$C=K_{AA}+K_{A\OO}x$ and
$(\Id-K_{\OO\OO})x=K_{\OO A}+v(1-C)$.
Substitution gives $(\Id-L)x=g_0$. Analytic continuation proves \eqref{tdc:eq:Cresolvent} wherever $\Id-L$ is invertible.

Solve the full block system for $\Id-L$ on $|q|=R$. If the right side has row norm at most one, \eqref{tdc:eq:schur} gives
\begin{align*}
 \norm{x_{\mathrm{core}}}_\infty
 &\le\frac{\kappa}{1-\kappa\delta}
 \left(1+\frac b{1-\tau}\right)
 =\frac{9063000}{17297}<525,\\
 \norm{x_{\mathrm{tail}}}_\infty
 &\le\frac{1+d(9063000/17297)}{1-\tau}
 =\frac{1483000}{121079}<525,
\end{align*}
where the equalities use the conservative rational bounds
$b=27/1000$, $d=21/1000$, $\tau=1/50$ and \eqref{tdc:eq:kappadelta}.
The bounded inverse constructed in this way agrees on square-summable vectors with the Hilbert-space inverse as in Section~\ref{tdc:sec:global}.

At $|z|\le4/5$, direct polynomial coefficient majorants give
\begin{equation}\label{tdc:eq:Cingredients}
 |K_{AA}|\le\frac{656}{625}<1.05,\quad
 \norm{K_{A\OO}}_1\le\frac{5386644}{1953125}<2.76,\quad
 \norm{g_0}_\infty\le\frac{123968}{78125}<1.59.
\end{equation}
The last calculation is finite: a state reaching $A$ in three steps has height at most four, and the other contribution is supported at $E$. Equations~\eqref{tdc:eq:Cresolvent}--\eqref{tdc:eq:Cingredients} imply
\[
 |C(q)|<1.05+2.76\cdot525\cdot1.59<2400\qquad(|q|=R).
\]
On the same circle, \eqref{tdc:eq:rho}--\eqref{tdc:eq:c} give
\[
 \left|\frac c{1-q/\rho}\right|
 \le\frac{c\rho}{R-\rho}
 <\frac{(0.184)(0.619)}{0.63-0.619}
 =\frac{14237}{1375}<10.355.
\]
The function $G(q)=C(q)-c/(1-q/\rho)$ is holomorphic on a neighborhood of the closed disk: the only pole inside was removed, and the boundary is zero-free for $N$. Its boundary modulus is below $2500$. Cauchy's coefficient estimate yields
$|[q^n]G|\le2500R^{-n}$ for every $n\ge0$, completing Theorem~\ref{tdc:thm:main}.

The error constant was chosen for transparency rather than sharpness. It is a global rigorous envelope, not a claim of a high-accuracy numerical approximation at small $n$.

\section{Every fixed finite spectral order}
\label{tdc:sec:spectral}
Meromorphic continuation contains more information than the first exponential term. The correct higher orders are contributions from further poles, potentially multiplied by polynomials and oscillations. An arbitrary series in $1/n$ would obscure this structure.

\begin{theorem}[Finite spectral expansion]\label{tdc:thm:spectral}
Fix $0<r<1$ such that $N$ has no zero on $|q|=r$. Let $\PP_r$ consist of the nonremovable poles of $C$ in that circle. For $\zeta\in\PP_r$, set
\[
 k_\zeta=\ord_\zeta N,\quad d_\zeta=\ord_\zeta D,
 \quad m_\zeta=k_\zeta-d_\zeta>0,
 \quad H_\zeta(q)=(1-q/\zeta)^{m_\zeta}C(q).
\]
Here $\ord_\zeta D=0$ when $D(\zeta)\ne0$. Define
\begin{equation}\label{tdc:eq:bcoeff}
 b_{\zeta,j}=\frac{(-\zeta)^{m_\zeta-j}}{(m_\zeta-j)!}
 H_\zeta^{(m_\zeta-j)}(\zeta),\qquad1\le j\le m_\zeta.
\end{equation}
There is a finite constant $M_r$ such that, for every $n\ge0$,
\begin{equation}\label{tdc:eq:spectral}
 a_n=\sum_{\zeta\in\PP_r}\sum_{j=1}^{m_\zeta}
 b_{\zeta,j}\binom{n+j-1}{j-1}\zeta^{-n}+E_{r,n},
 \qquad |E_{r,n}|\le M_r r^{-n}.
\end{equation}
One may take $M_r=\max_{|q|=r}|H_r(q)|$, where
\[
 H_r(q)=C(q)-\sum_{\zeta\in\PP_r}\sum_{j=1}^{m_\zeta}
 b_{\zeta,j}(1-q/\zeta)^{-j}.
\]
For $r=R$, the sole term is $c\rho^{-n}$ and $M_R=2500$ is valid.
\end{theorem}
\begin{proof}
The function $N$ is holomorphic in the unit disk, not identically zero because $N(0)=1$. Its zeros in the fixed closed disk are finite. At each zero, the quotient $D/N$ has exactly the surviving pole order $k_\zeta-d_\zeta$ when positive, and otherwise is removable. Taylor-expanding the analytic function $H_\zeta$ about $\zeta$ yields \eqref{tdc:eq:bcoeff}. Subtracting these finite principal parts removes every singularity in a neighborhood of the closed disk. The elementary coefficient identity
\[
 [q^n](1-q/\zeta)^{-j}=\binom{n+j-1}{j-1}\zeta^{-n}
\]
and Cauchy's estimate for $H_r$ prove \eqref{tdc:eq:spectral}.
\end{proof}

This is an all-orders statement at every \emph{fixed spectral accuracy}: select a fixed zero-free radius below one and include all its surviving poles. It allows repeated poles, conjugate pairs, and numerator cancellations. It neither says that an infinite pole sum converges nor gives uniform estimates as $r\uparrow1$. Only the dominant pole and $M_R$ have been numerically certified here. No simplicity assumption is made for later zeros.

For the single leading carrier, \eqref{tdc:eq:error} implies
$a_n/(c\rho^{-n})=1+o(n^{-J})$ for every fixed $J$. Thus every nonconstant coefficient of any relative algebraic $1/n$ expansion around that carrier is zero. The nontrivial higher information is the additional exponential spectral terms in \eqref{tdc:eq:spectral}.

\section{Strict growth and threshold inversion}
\label{tdc:sec:inverse}
\subsection{The two empirical statements and the effective growth gate}
For the exact rational $x=0.618034$,
\[
 x^2+x-1=\frac{6289}{250000000000}>0.
\]
Since the positive zero of $x^2+x-1$ is $1/\phi$, this proves
$1/\phi<x<\rhol<\rho$, and hence $\lambda<\phi$. Equation~\eqref{tdc:eq:error} gives $a_n/\phi^n\to0$ and
$a_{n-j}/a_n\to\rho^j$ for each fixed $j$. The signed residual therefore tends to $1-\rho-\rho^2$. This function decreases on the positive real line; evaluating its endpoints with exact rational arithmetic gives \eqref{tdc:eq:empirical}. In particular $(a_{n-1}+a_{n-2})/a_n$ tends to $\rho+\rho^2\ne1$.

To make eventual positivity and growth effective, set
\[
 Q_n=\frac M\cl\left(\frac{\rhoh}{R}\right)^n.
\]
Exact rational comparisons at $n=600$ give
\begin{equation}\label{tdc:eq:growthgate}
 Q_{600}<0.138<1,\qquad
 Q_{600}(1+R^{-1})<0.356406<\rhoh^{-1}-1.
\end{equation}
The left sides decrease with $n$. Equation~\eqref{tdc:eq:error} consequently gives $a_n>0$ and
\begin{align*}
 a_{n+1}-a_n
 &\ge c\rho^{-n}(\rho^{-1}-1)-MR^{-n}(1+R^{-1})>0
 \qquad(n\ge600).
\end{align*}
This proves Corollary~\ref{tdc:cor:empirical} without assuming positivity at any uncomputed earlier index.

Define the exact integer
\begin{equation}\label{tdc:eq:Y0}
 Y_0=\left\lceil\ch\rhol^{-600}+MR^{-600}\right\rceil.
\end{equation}
For clarity, its decimal expansion is the concatenation of the following five successive blocks of $25$ digits:
\begin{center}\ttfamily
5135418644989400288698135\\
4167987143108411082531436\\
8123350447210340481936646\\
4412227270131795726763531\\
3735694718693217016233740
\end{center}
For every $0\le n\le600$, the upper bound
$\ch\rhol^{-n}+MR^{-n}$ is at most $Y_0$. Thus every such $a_n$ is below any threshold $y>Y_0$. This large conservative threshold is a guaranteed prefix exclusion, not a numerical estimate of the optimal onset.

\subsection{A pointwise inverse with two ceilings}
Put $\lambda=\rho^{-1}$ and $L=\log\lambda$. For $y>Y_0$ let
\[
 n(y)=\min\{n\ge0:a_n\ge y\}.
\]
It exists because the sequence is eventually strictly increasing and unbounded.

\begin{theorem}[Explicit threshold enclosure]\label{tdc:thm:inverse}
For each $y>Y_0$, there are unique roots $x_-(y),x_+(y)>600$ of
\begin{align}
 c\lambda^{x_-}+MR^{-x_-}&=y,\label{tdc:eq:upperroot}\\
 c\lambda^{x_+}-MR^{-x_+}&=y.\label{tdc:eq:lowerroot}
\end{align}
They satisfy the pointwise integer enclosure
\begin{equation}\label{tdc:eq:twoceil}
 \left\lceil x_-(y)\right\rceil\le n(y)\le
 \left\lceil x_+(y)\right\rceil.
\end{equation}
The constants $c,\rho$ in these equations are their exact determinant-defined values; their interval enclosures are available for certified numerical evaluation.
\end{theorem}
\begin{proof}
The plus envelope is strictly increasing. For real $x\ge600$, the relative subtracted term in the minus envelope is
$(M/c)(\rho/R)^x\le Q_{600}<1$ and decreases with $x$. Also
$0<\log(1/R)<L$. The minus envelope is positive and its derivative
\[
 c\lambda^x L-MR^{-x}\log(1/R)>0.
\]
Both envelopes tend to infinity and have value below $y$ at $600$, by \eqref{tdc:eq:Y0}. Thus their stated roots exist and are unique. Any integer $n>600$ below $x_-$ has upper envelope below $y$, so $a_n<y$. All earlier indices were already excluded. At $n=\lceil x_+\rceil$, the lower envelope is at least $y$, so $a_n\ge y$. This proves \eqref{tdc:eq:twoceil}, including the case when a root is exactly an integer.
\end{proof}

A single ceiling around an expression with an unspecified $O$-term would not give a pointwise result near integers. Keeping the two ceilings is essential. Even when the horizontal interval is short, it can straddle an integer.

\subsection{All orders for the two bounding roots}
Define
\begin{equation}\label{tdc:eq:inverseparams}
 x_0=\frac{\log y-\log c}{L},\quad
 \alpha=\frac{\log(1/R)}{L}\in(0,1),\quad
 \beta=1-\alpha=\frac{\log(R/\rho)}{L}>0,\quad
 \mu=Mc^{-\alpha}y^{-\beta}.
\end{equation}
Generalized binomial coefficients have their polynomial meaning,
$\binom{v}{j}=v(v-1)\cdots(v-j+1)/j!$ for $j\ge1$, and value one for $j=0$.

\begin{proposition}[Lagrange expansions]\label{tdc:prop:lagrange}
For every fixed integer $K\ge1$, as $y\to\infty$,
\begin{align}
 x_-(y)&=x_0+\frac1L\sum_{k=1}^K
 \frac{(-1)^k}{k}\binom{\alpha k-1}{k-1}\mu^k
 +O_K(\mu^{K+1}),\label{tdc:eq:minusseries}\\
 x_+(y)&=x_0+\frac1L\sum_{k=1}^K
 \frac1k\binom{\alpha k-1}{k-1}\mu^k
 +O_K(\mu^{K+1}).\label{tdc:eq:plusseries}
\end{align}
These series converge for sufficiently small $|\mu|$. Their first terms are
\[
 x_-=x_0-\frac\mu L+\frac{\alpha-1/2}{L}\mu^2+O(\mu^3),\qquad
 x_+=x_0+\frac\mu L+\frac{\alpha-1/2}{L}\mu^2+O(\mu^3).
\]
\end{proposition}
\begin{proof}
Set $v=\exp(L(x-x_0))$. The equations become $v+\mu v^\alpha=1$ and $v-\mu v^\alpha=1$. Near $v=1$, choose the analytic power defined by the logarithm with $\log1=0$. Write $w=v-1$; then
\[
 w=\epsilon\mu(1+w)^\alpha,
 \quad\epsilon=-1\ \text{for }x_-,\quad\epsilon=1\ \text{for }x_+.
\]
The implicit derivative with respect to $w$ is one at $(w,\mu)=(0,0)$, so the solution and its logarithm are analytic near zero. Lagrange inversion for $\log(1+w)$ gives
\[
 [\mu^k]\log(1+w)=\frac{\epsilon^k}{k}
 [w^{k-1}](1+w)^{\alpha k-1}
 =\frac{\epsilon^k}{k}\binom{\alpha k-1}{k-1}.
\]
Multiplying by $1/L$ proves the formulas and their fixed-order analytic remainder bounds.
\end{proof}

These are expansions of exact \emph{bounding roots}, not asymptotic equalities for the integer-valued step function $n(y)$. Adding terms to either Lagrange series improves its evaluation but does not reduce the sequence's underlying $O(y^{-\beta})$ horizontal uncertainty. Improving that width requires a sharper coefficient error or further spectral information.

\subsection{Inversion at each fixed spectral accuracy}
Fix $\rho<r<1$ as in Theorem~\ref{tdc:thm:spectral}. Its finite spectral carrier has a real smooth interpolation $A_r(t)$ that must be specified, since arbitrary complex powers would otherwise be ambiguous. Put
\[
 P_j(t)=\binom{t+j-1}{j-1}.
\]
For a positive real pole use $b_{\zeta,j}P_j(t)\zeta^{-t}$. For a negative real pole use
$b_{\zeta,j}P_j(t)|\zeta|^{-t}\cos(\pi t)$. For each conjugate pair choose the upper-half-plane pole $\zeta=|\zeta|e^{i\theta}$ with $0<\theta<\pi$ and use
\[
 2\Re\bigl(b_{\zeta,j}P_j(t)|\zeta|^{-t}e^{-i\theta t}\bigr).
\]
Summing these terms defines $A_r(t)$ and agrees exactly with the carrier in \eqref{tdc:eq:spectral} at integers. Real coefficients of $C$ ensure the required conjugacy of Laurent coefficients.

Every nonleading pole has modulus greater than $\rho$. Its interpolated term and derivative are therefore exponentially smaller than $c\lambda^t$, even after fixed polynomial factors. Hence
\[
 A_r(t)\sim c\lambda^t,\qquad A_r'(t)\sim cL\lambda^t.
\]
The two functions $A_r(t)\pm M_r r^{-t}$ are eventually strictly increasing and unbounded. For sufficiently large $y$, their respective large roots $x_{r,-},x_{r,+}$ give
\begin{equation}\label{tdc:eq:spectralinverse}
 \lceil x_{r,-}(y)\rceil\le n(y)\le\lceil x_{r,+}(y)\rceil.
\end{equation}
One excludes the finitely many earlier indices before applying the same argument as for \eqref{tdc:eq:twoceil}. This exclusion and the finer constants are existential here, rather than additional certified numerical tables.

Both roots equal $x_0+o(1)$. The derivative of the smooth carrier near them is asymptotic to $Ly$, whereas the vertical separation of the two envelopes is $O(r^{-x_0})$. The mean value theorem gives
\begin{equation}\label{tdc:eq:betar}
 x_{r,+}(y)-x_{r,-}(y)=O(y^{-\beta_r}),\qquad
 \beta_r=\frac{\log(r/\rho)}{\log(1/\rho)}\in(0,1).
\end{equation}
There are zero-free radii approaching one because zeros are discrete in the open disk, and along such radii $\beta_r\uparrow1$. Meromorphic continuation confined to $|q|<1$ does not supply arbitrary inverse errors $O(y^{-K})$ for $K>1$. Further implicit expansions may use the explicit finite carrier; at no point is the discrete unknown remainder $E_{r,n}$ differentiated.

\begin{remark}[{[write]} The inverses against the transseries volume, 7~October 2026]\label{tdc:rem:transseries}
The volume is
\path{Analysis/Transseries/docs/series-and-transseries/Transseries_And_Inversion/transseries_and_inversion.tex};
its symbols carry the subscript ``vol''. Statement by statement:
\begin{enumerate}[label=(\alph*)]
\item \emph{Instance.} For $y>Y_0$, $n(y)$ is the integer staircase of
 \file{p0:def:three-inverses}(1) with $A_n=a_n$ and
 $n_{1,\mathrm{vol}}=600$: by Corollary~\ref{tdc:cor:empirical} the sequence
 is strictly increasing from $n=600$, and by~\eqref{tdc:eq:Y0} no earlier
 index reaches $y$. So \file{p0:thm:staircase}(1) gives
 $n(y)=\lceil\nu_{\mathcal A}(y)\rceil$ for every admissible interpolation
 $\mathcal A_{\mathrm{vol}}$.
\item \emph{Analogue.} Theorem~\ref{tdc:thm:inverse} is a direct two-ceiling
 bracket at the integers about the roots of two explicit envelopes, which do
 not interpolate the sequence: an analogue of \file{p0:thm:staircase}(2), not
 an instance, with fully explicit constants. The same holds
 for~\eqref{tdc:eq:spectralinverse}, whose carrier $A_r$ interpolates the
 spectral carrier, not $a_n$.
\item \emph{Instances.} The centre $x_0$ of~\eqref{tdc:eq:inverseparams} is the
 case $b_{\mathrm{vol}}=0$ of \file{p0:thm:lambert-core}
 ($a_{\mathrm{vol}}=L$, target $\log y-\log c$). In the proof of
 Proposition~\ref{tdc:prop:lagrange}, the equation $w=\epsilon\mu(1+w)^\alpha$
 is a formal instance of \file{p0:thm:core-reversion} with
 $\Lambda_{\mathrm{vol}}=1$, $h_{\mathrm{vol}}=0$,
 $f_{\mathrm{vol}}(t,u)=-\epsilon t(1+u)^\alpha$ and $t=\mu$ (over a ring
 containing $\alpha$); the source's coefficient formula is the
 Lagrange--B\"urmann formula for $\log(1+w)$, and its convergence (analytic
 implicit function theorem) is the source's own.
\end{enumerate}
No novelty is claimed for any of these inversions. The staircase arithmetic
is formalized generically (for an arbitrary strictly monotone function) as
\file{Fabius.staircase_ceil} in
\path{Analysis/FabiusFunction/Lean/FabiusFunction/StaircaseInversion.lean};
nothing about $a_n$ is formalized.
\end{remark}

\section{A local analytic span marker}
\label{tdc:sec:marker}
The exact span variable from Section~\ref{tdc:sec:model} can be retained. It leads to a signed analytic extension, independently of any probability interpretation. A primitive size-$m$ path has $m-1$ compressed transitions, so
\[
 J(t,q)=t e_A^{\mathsf T}(\Id-tK(z))^{-1}h(z),\qquad q=z^2.
\]
Cramer's replacement column is still $h=e_A+z^3e_E$. Its row elimination now gives
$\Id-tK_{\OO\OO}+tz^3e_EK_{A\OO}=\Id-tL$. Thus
\begin{equation}\label{tdc:eq:markeridentity}
 C(t,q)=1-\frac{D(t,q)}{N(t,q)},\quad
 D(t,q)=\det(\Id-tK(\sqrt q)),\quad
 N(t,q)=\det(\Id-tL(\sqrt q)).
\end{equation}
These determinants are jointly holomorphic for $t\in\mathbb C$, $|q|<1$: the trace-norm series converges uniformly on compact marker sets, and parity removes the square-root branch as before. Their values at $q=0$ are one. For each fixed bounded marker set, \eqref{tdc:eq:markeridentity} agrees with the source formal series in a sufficiently small common $q$ neighborhood.

\begin{theorem}[Local marker continuation]\label{tdc:thm:marker}
There is an $\epsilon>0$ and analytic functions $\rho(t),c(t)$ for $|t-1|<\epsilon$, extending $\rho,c$, such that $N(t,q)$ has exactly one zero in $|q|<R$, simple and equal to $\rho(t)$, and $D(t,\rho(t))\ne0$. With
\[
 c(t)=\frac{D(t,\rho(t))}{\rho(t)\,\partial_q N(t,\rho(t))},
 \qquad A_n(t)=[q^n]C(t,q),
\]
one has, on every sufficiently small closed marker disk,
\begin{equation}\label{tdc:eq:markerexpansion}
 A_n(t)=c(t)\rho(t)^{-n}+O(R^{-n})
\end{equation}
uniformly in $t$. Each fixed marker derivative has the derivative of the displayed leading carrier plus an $O(R^{-n})$ error, uniformly on any still smaller closed marker disk.
\end{theorem}
\begin{proof}
At $t=1$ the denominator has a simple zero at $\rho$ and no zero on $|q|=R$. The analytic implicit-function theorem gives a local simple-zero branch. Uniform continuity on the compact boundary preserves its nonvanishing after shrinking the marker disk; the argument principle preserves the count one. The numerator stays nonzero at this branch by continuity. After a further shrink, $\rho(t)$ stays away from both zero and the boundary and $c(t)$ stays nonzero.

Subtract $c(t)/(1-q/\rho(t))$ from \eqref{tdc:eq:markeridentity}. Its apparent singularity on the simple-zero branch is jointly removable. The resulting function is jointly holomorphic inside the circle and uniformly bounded on its boundary over a closed smaller marker disk. Cauchy's coefficient estimate proves \eqref{tdc:eq:markerexpansion}. For derivatives, first establish that bound on an intermediate closed marker disk and apply Cauchy's formula in $t$ on circles contained in it. This proves the derivative bounds on a still smaller disk with derivative-dependent constants.
\end{proof}

No numerical value of $\epsilon$, marker derivatives, or variance is certified. The support lemma makes each $A_n(t)$ a polynomial, but no general coefficientwise positivity, probability law, central limit theorem, or probabilistic nondegeneracy follows from this signed determinant argument alone.

\section{Reproducibility and the boundary of the claims}
\label{tdc:sec:repro}
\subsection{What the accompanying code proves}
The accompanying archive contains this \LaTeX{} source, the PDF, exact finite polynomial and adjugate data, frozen dyadic numerical proposals, independently written coefficient generation, and all acceptance checks. The code reconstructs the operators from their stated transition rules; it does not execute or redistribute the third-party jOEIS implementation. No third-party paper is included.

The mandatory proof replay uses Python's standard library. Its global portion uses only integers and rational numbers. The local portion uses outward dyadic intervals with integer endpoints and residual-tested exact proposals. Optional \texttt{mpmath} code can regenerate proposals, but a wrong or corrupted proposal cannot pass merely because it came from that numerical library: the integer residual tests are the acceptance gate. No mandatory symbolic algebra system or interval-trigonometric library is required.

The main reproducibility command is
\begin{center}\texttt{python3 reproduce.py}\end{center}
from the extracted archive directory. It runs the full certificate suite with Python site initialization disabled, in both normal and optimized modes, then compares deterministic output bytes. The package README gives the individual commands, output inventory, environment details, and optional proposal-generation route. Each failing acceptance condition raises an explicit exception; none relies on Python's removable \texttt{assert} statement.

The key checks are as follows.
\begin{enumerate}
\item Reconstruct the complete height-eight state list and signed operator, verify the adjugate identity, and prove the supplied polynomial is the genuine determinant using its degree bound and $240$ exact evaluations. A separate rational Gaussian implementation verifies those evaluations; the polynomial can also be regenerated exactly.
\item Evaluate every rational-circle node and its derivative majorant, count the signed polygon crossings including closure, and verify every tail and Schur inequality. Independent winding checks use multiple rays.
\item Verify the local determinant signs and nonsingular numerator, inverse residuals, solution residuals, derivative tail bounds, and widened amplitude interval.
\item Verify the boundary resolvent estimate, $M=2500$, the exact golden-ratio separation, residual limit bounds, strict-growth inequalities, and the full integer $Y_0$.
\item Generate $65$ exact cluster coefficients independently from the deficit paths. The first $40$ agree with the displayed OEIS terms inspected in the source screen; terms beyond $40$ are generated values, not a claimed b-file comparison.
\end{enumerate}
The finite coefficient check is supplementary. The proof of source identity in Section~\ref{tdc:sec:model} is what identifies the model at every index.

Positive arithmetic tests compare independent algorithms. Negative tests alter states, coefficients, adjugate entries, common polynomial factors, circle nodes, crossings, local proposals, interval conditions, and coefficient data. In particular, multiplying both the adjugate and its proposed denominator by $1+z^2$ preserves \eqref{tdc:eq:adjugate} but must fail the genuine-determinant test. A closing-chord test is accurately described as deleting a vertex so that the new chord crosses the origin; actual loop closure is also exercised by the independent negative-ray winding calculation. Tests run under both normal Python and \texttt{-O}.

Hashes bind files and help detect accidental change. They are not mathematical evidence in place of the checks. Stored \texttt{PASS} outputs are convenient records, but the reproducibility command recomputes them rather than treating them as axioms. Deterministic builds and negative tests reduce implementation risk; they do not claim a machine-checked proof in a formal theorem prover.

\subsection{Source identification and bounded comparison}
\label{tdc:sec:sources}
The OEIS entry~\cite{oeis} identifies the row-word inversion cluster sequence, its hard-rod recursion, the displayed prefix, and empirical golden-ratio and Fibonacci comparisons. The golden-ratio form was described there as checked to four significant figures over indices $30$--$40$. Its separate observation of a small relative Fibonacci residual is quantitatively consistent with \eqref{tdc:eq:empirical}, despite the failure of a strict equivalent.

The original Python attachment linked from the record was not accessible through the attempted permitted route. The entire jOEIS translation~\cite{joeis}, explicitly crediting Brydensholt's Python, was inspected. The inspected content has Git blob identifier
\begin{center}\texttt{e59ad7a42e0c96f832b97b6457ea9b5cb1d4a851}.
\end{center}
Its full source is not bundled. The live OEIS record was separately inspected in a browser; a cached primary-source search rendition was also checked. The source screen was performed on 3 October 2026. The linked $1$--$100$ b-file was not independently used to verify all those terms.

\begin{writenote}[Added 7 October 2026]
The write fetched the b-file (Irvine, $1\le n\le100$) on 7~October 2026: all
$100$ terms agree with its own exact generator, and terms $1$--$65$ with the
shipped \file{data/exact_coefficients.json}
(Remark~\ref{tdc:rem:refuted}(d)).
\end{writenote}

Targeted searches of the exact sequence identifier, author, row-word inversion, and defect-cluster terminology did not locate a directly covering asymptotic theorem in the sources checked. Related OEIS entries A393486 and A393487 concern the surrounding span/weight model. A seemingly related cluster-algebra report concerns a different annular cluster-variable problem~\cite{clustercomparison}. An inspected 2026 Fibonacci ribbon-tableau manuscript uses a different shape, ballot condition, and enumerated statistic~\cite{ribboncomparison}. Its opening theorem/model comparison does not amount to a full proof audit. These bounded negative and model-comparison results do not establish worldwide novelty or exhaustive duplicate clearance.

For the classical determinant background, relevant portions of Bornemann's primary preprint, Sections 2--4, were inspected~\cite{bornemann}. The general trace-class facts are sufficient here. Its separate erratum corrects a specialized positive self-adjoint lemma~\cite{bornemannerratum}; that lemma is not applied to $L$. Simon's 1977 paper~\cite{simon} is a classical reference, but only its publisher abstract was available in this source screen, so no unverified theorem-number attribution is made to it.

\subsection{Useful directions left open}\label{tdc:sec:open}
\begin{enumerate}
\item \textbf{Sharper certified numerics.} Larger cores or optimized rational contours could narrow the pole and amplitude brackets, reduce $M$, and bring the guaranteed monotonicity and inverse thresholds closer to observed behavior. This requires new global and tail certificates, rather than merely a more precise floating root.
\item \textbf{Subdominant poles.} Locate and certify the next nonremovable poles, their orders, and residues. Theorem~\ref{tdc:thm:spectral} then turns them directly into useful oscillatory corrections and tighter inverse enclosures. Numerator cancellation must be tested at each candidate.
\item \textbf{Finite-index positivity.} Determine whether every source cluster coefficient is nonnegative, and whether it admits a positive combinatorial interpretation. The signed rank-one formulation proves neither assertion. Eventual positivity of the univariate coefficients already follows from the present estimate.
\item \textbf{Span statistics.} Certify a quantitative marker disk and marker derivatives. A probability or central limit theorem additionally requires a positive model and suitable nondegeneracy, beyond Theorem~\ref{tdc:thm:marker}.
\item \textbf{The unit circle.} Determine whether continuation through any part of $|q|=1$ is possible, whether poles accumulate there, or whether a natural boundary occurs. The present trace-class estimate only applies strictly inside that circle.
\item \textbf{Exact structure.} Seek further identities for the determinant quotient or its growth constants, without assuming an algebraic value suggested by numerical proximity. The strict inequality $\lambda<\phi$ illustrates the importance of separating a near match from an exact asymptotic law.
\end{enumerate}

\begin{writenote}[Added 7 October 2026, under Vladimir's standing rule of 4 October 2026]
The six directions above are the source's open questions, credited to it and
kept as stated. The write found no unproved claim of the source outside them
and no wrong one; the two OEIS formula lines are refuted in
Remark~\ref{tdc:rem:refuted}. It adds two finite questions that its own
computations leave open.
\begin{enumerate}[label=\arabic*.,start=7]
\item \emph{Monotonicity below $600$.} The write's exact terms show
 $a_{n+1}>a_n$ for $4\le n<300$ (and $a_1<a_2>a_3=a_4$);
 Corollary~\ref{tdc:cor:empirical} gives it for $n\ge600$. Missing: the exact
 terms for $300\le n\le600$ (a longer run of the same generator), which would
 give $a_{n+1}>a_n$ for every $n\ge4$.
\item \emph{The OEIS residual comment for every $n\ge30$.} It holds for
 $30\le n\le300$ by exact computation and for all large $n$ by the limit; the
 explicit envelope~\eqref{tdc:eq:error} makes it effective only for $n$
 near $1000$ or beyond. Missing: exact terms up to that range, or a sharper
 envelope (direction~1).
\end{enumerate}
\end{writenote}

\begin{writenote}[Added after the independent check, 7 October 2026]
Questions~7 and~8 are settled, both resting on the computer certificate of
Theorem~\ref{tdc:thm:main} for their tails. The check computed
$a_1,\dots,a_{1250}$ exactly with its own weight-ordered generator (the
walk~\eqref{tdc:eq:edgesraw} without a step count or height truncation, states
pruned only by the gauge bound $2W\ge g(u)+1$ on the deficit still needed to
return; it equals the b-file in all $100$ terms and the write's terms). It
finds $a_{n+1}>a_n$ for $4\le n<1250$, which with
Corollary~\ref{tdc:cor:empirical} gives $a_{n+1}>a_n$ for every $n\ge4$
(Question~7); and $|a_n-a_{n-1}-a_{n-2}|/a_n<0.0002$ for $30\le n\le1250$,
which with the envelope bound for $n\ge982$ given in
Remark~\ref{tdc:rem:refuted}(c) proves the OEIS comment's inequality for
every $n\ge30$ (Question~8).
\end{writenote}

\begin{writenote}[Independent check of the write, 7 October 2026]
An adversarial check made by the intake after the write (commit
\file{5d3107954}) re-derived the statements marked [write] with its own code,
after fetching again A394326 (revision~\#15) and its b-file, A393486
(\#28) and A393487 (\#34). \emph{Remark~\ref{tdc:rem:oeis}.} Every
quotation, revision, date and author line confirmed. \emph{The refutation}
(Remark~\ref{tdc:rem:refuted}), checked hardest. The certificate chain
replayed from the shipped files (Route~B of the README: global, local,
remainder, inverse, coefficients, independent global check, local tests and
the regeneration of $N_8$, all equal to the delivered records, \texttt{PASS},
49~s). Independently of the certificate: the exceptional row
\eqref{tdc:eq:exceptional} re-derived symbolically; the zero of
height-truncated $N$ computed at heights $14$, $20$, $26$ by floating
determinants, $0.6180397634035489691315\ldots$ each time (with
$N_H(1/\phi)>0$ and $D_{20}=-0.43190569\ldots$ there, inside the source's
interval); and the ratio $a_{n-1}/a_n$ of the exact terms equal to the same
value to $60$ digits at $n=1100$ (consecutive errors falling by about
$10^{-9}$ per $100$ indices, so the next pole has modulus near $0.755$). This
is corroboration, not proof; it sits inside the certified
bracket~\eqref{tdc:eq:rho}, $5.77\ldots\cdot10^{-6}$ above $1/\phi$. The proof of~(a) and~(b), the
exact rationals ($6289/(25\cdot10^{10})$, \eqref{tdc:eq:lambda}, the band
of~\eqref{tdc:eq:empirical}, \eqref{tdc:eq:growthgate}, all $125$ digits of
$Y_0$, and $a_{601}$ the first term above $Y_0$), $\lambda/\phi<0.9999907$,
$(\lambda/\phi)^{10}>0.99990$ and every exact value printed in~(c) and~(d)
confirmed (truncations). The ``(empirical)'' Fibonacci line is refuted in
its natural reading, as stated. Two additions, by dated notes: the comment's
inequality and the decrease of $a_n/\phi^n$ for every $n\ge30$
(Remark~\ref{tdc:rem:refuted}(c)), and Questions~7 and~8 settled (above).
\emph{Remark~\ref{tdc:rem:transseries}.} (a)--(c) re-derived against the
volume: $n_{1,\mathrm{vol}}=600$ is admissible ($a_n<y$ for $n\le600$ by the
definition of $Y_0$); \file{p0:thm:lambert-core}(1) with $b=0$; and
\file{p0:thm:core-reversion} with $\Lambda=1$, $h=0$, $f=-\epsilon
t(1+u)^\alpha$ is a formal instance as stated. \emph{Provenance.} Archive
facts, the staged files (byte-identical), the 71 delivered label numbers and
59 references (aux of a build of the delivered text, 21 pages) and the
comparison report's blob confirmed. One inconsistency, recorded by a dated
note in Section~\ref{tdc:sec:provenance}: the sentence that the write did not
rerun the certificate contradicts the README and the write's commit message.
No error was found in the source's proofs.
\end{writenote}

\begin{thebibliography}{9}
\raggedright
\bibitem{oeis}
Morten Brydensholt, \emph{A394326}, sequence record on connected tableau defect-cluster types, The On-Line Encyclopedia of Integer Sequences.
\url{https://oeis.org/A394326}. Source screen: 3 October 2026 (read by the write on 7~October 2026, revision \#15, with its b-file).

\bibitem{joeis}
Sean A. Irvine, \emph{A394326.java}, jOEIS, based on Morten Brydensholt's Python program.
\url{https://github.com/archmageirvine/joeis/blob/master/src/irvine/oeis/a394/A394326.java}.
Inspected Git blob: \texttt{e59ad7a42e0c96f832b97b6457ea9b5cb1d4a851}.

\bibitem{bornemann}
Folkmar Bornemann, \emph{On the Numerical Evaluation of Fredholm Determinants}, Mathematics of Computation \textbf{79} (2010), 871--915. Primary preprint:
\url{https://arxiv.org/abs/0804.2543}; PDF:
\url{https://arxiv.org/pdf/0804.2543}.

\bibitem{bornemannerratum}
Folkmar Bornemann, \emph{Erratum to On the Numerical Evaluation of Fredholm Determinants}.
\url{https://www-m3.ma.tum.de/bornemann/errata/FredDetErratum.pdf}.

\bibitem{simon}
Barry Simon, \emph{Notes on infinite determinants of Hilbert space operators}, Advances in Mathematics \textbf{24} (1977), 244--273.
\url{https://doi.org/10.1016/0001-8708(77)90057-3}.

\bibitem{clustercomparison}
\emph{Cluster-variable log-concavity counterexample}, ProveIt repository, comparison source concerning a different cluster-algebra model.
\url{https://github.com/VladimirReshetnikov/ProveIt/blob/main/SetTheory/Cardinals/docs/reports/log-concavity-and-unimodality/cluster-variable-log-concavity-counterexample/article.tex}.
Opening 120 lines inspected; content blob \texttt{4d75452584d2b7b2bd01b5c565b47a31578611cb}.

\bibitem{ribboncomparison}
Alex Chengyu Li, \emph{Exact Enumeration and Fixed-Rank Asymptotics of Ballot-Admissible Fibonacci Ribbon Tableaux}, 2026 manuscript, comparison source.
\url{https://github.com/crabsatellite/fibonacci-ribbon-ballot-enumeration}.
README and opening 190 manuscript lines inspected; manuscript content blob
\texttt{22c6ec3c863c2ab698630f9da817a039d89b0b67}.
\end{thebibliography}
\end{document}
